\section{Implementacja analizy wyników}

W podanym rozdziale opisujemy końcowy etap implementacji, wspólny dla wszystkich implementowanych rozwiązań metod uczenia zespołowego, obejmujący ewaluację modeli oraz analizę i wizualizację uzyskanych wyników. Ten proces został zaimplementowany w dwóch głównych etapach i koncentruje się na porównaniu skuteczności klasyfikacji dla metod naukowych, jak i autorskich. Proces ten wykorzystuje dane z kluczowych zestawów danych do metryk wydajności oraz do krzywych precyzji\-odwołania, organizując wyniki według kategorii zestawów danych dla standaryzowanej oceny.

Dane wejściowe do analizy wyników są generowane przez każdą z 20 metod (\texttt{metoda1} do \texttt{metoda20}), które zwracają wytrenowany model oraz dane testowe w następującej postaci:
\begin{itemize}
	\item \textbf{model} – wytrenowany klasyfikator, reprezentujący implementację danej metody, gotowy do generowania predykcji.
	\item \textbf{X\_test} – tablica zawierająca cechy zbioru testowego, które mogą być wektorami TF-IDF lub osadzeniami kontekstowymi z BERT, RoBERTa czy Sentence Transformers, wykorzystywanymi do oceny modelu.
	\item \textbf{y\_test} – tablica zawierająca etykiety binarne zbioru testowego: 0 dla fałszywych wiadomości, 1 dla prawdziwych; służące jako punkt odniesienia dla predykcji.
\end{itemize}
Te dane są przekazywane do dalszej ewaluacji i analizy, a ich spójna struktura umożliwia standaryzowane porównanie wyników.

Pierwszym krokiem jest ewaluacja modeli przy użyciu funkcji \texttt{evaluate\_model}, wywoływanej dla każdej metody po jej uruchomieniu. Poniżej znajduje się szczegółowy opis parametrów przekazywanych do funkcji \texttt{evaluate\_model}:
\begin{itemize}
	\item \textbf{model} – zwrócony parametr wyniku metody, szczegółowo opisaliśmy poprzednio.
	\item \textbf{X\_test} – zwrócony parametr wyniku metody, szczegółowo opisaliśmy poprzednio.
	\item \textbf{y\_test} – zwrócony parametr wyniku metody, szczegółowo opisaliśmy poprzednio.
	\item \textbf{run\_type} – identyfikator typu uruchomienia, używany do oznaczania wyników w plikach wyjściowych.
	\item \textbf{output\_path} – ścieżka do katalogu, w którym zapisywane są wyniki, umożliwiająca organizację danych wyjściowych.
	\item \textbf{start\_time} – znacznik czasu rozpoczęcia uruchomienia, służący do obliczania całkowitego czasu wykonania.
\end{itemize}

Wspominając o szczegółach tego kroku, funkcja rozpoczyna od przeprowadzenia walidacji krzyżowej, obliczając średnią dokładność (\texttt{cv\_accuracy\_mean}) oraz odchylenie standardowe (\texttt{cv\_accuracy\_std}). Walidacja krzyżowa to metoda oceny modelu, w której dane testowe dzieli się na 5 równych części – model trenuje się na 4 z nich, a testuje na jednej, powtarzając proces 5 razy, za każdym razem zmieniając testowaną część. Dzięki temu można ocenić, jak stabilny i niezawodny jest model w różnych warunkach, zamiast sprawdzać go tylko na jednym zestawie danych. Następnie model generuje predykcje prawdopodobieństw – dla standardowych struktur danych za pomocą \texttt{predict\_proba}, a dla zaawansowanych struktur, takich jak sieci neuronowe, za pomocą \texttt{predict}. Te prawdopodobieństwa są konwertowane na etykiety binarne przy progu 0.5, gdzie wartość powyżej progu oznacza klasę pozytywną (1). Obliczane są metryki oceny, w tym:

\begin{itemize}
	\item \textbf{Dokładność (Accuracy)} – ile przewidywań było poprawnych.
	\item \textbf{Precyzja (Precision)} – jak często, gdy model mówi "prawda", ma rację.
	\item \textbf{Recall} – jak dobrze znajduje prawdziwe wiadomości.
	\item \textbf{Wynik F1 (F1-Score)} – równowaga między precyzją a czułością.
	\item \textbf{Pole pod krzywą ROC-AUC (ROC-AUC)} – jak dobrze model odróżnia klasy.
	\item \textbf{Współczynnik korelacji Matthewsa (MCC)} – miara zgodności przewidywań z rzeczywistością.
	\item \textbf{Logarytmiczna strata (Log Loss)} – jak bardzo model się myli.
	\item \textbf{Współczynnik Kappa Cohena (Cohen's Kappa)} – jak dobrze model przewyższa losowe zgadywanie.
\end{itemize}
Dodatkowo mierzony jest czas wykonania (\texttt{Execution Time}), a wyniki walidacji krzyżowej są dołączane do metryk. Wszystkie wartości są zapisywane do pliku CSV w formacie \texttt{runX\_results.csv}, a krzywa Precision-Recall (\texttt{precision\_recall\_curve}) jest generowana i zapisywana jako \texttt{runX\_pr\_curve.csv}. Na koniec metryki są wyświetlane w konsoli, co pozwala na szybką weryfikację skuteczności modelu. W tym procesie \texttt{model} służy do generowania predykcji, \texttt{X\_test} i \texttt{y\_test} wspierają obliczanie metryk i walidację krzyżową, \texttt{run\_type} oraz \texttt{output\_path} określają format i lokalizację zapisu wyników, a \texttt{start\_time} pozwala zmierzyć czas wykonania. Kluczowe parametry wykorzystane w procesie ewaluacji modeli obejmują:
\begin{itemize}
	\item \textbf{n\_splits=5} – liczba podziałów w walidacji krzyżowej, określająca, ile razy dane są dzielone na zestawy treningowe i testowe w celu obliczenia średniej dokładności.
	\item \textbf{random\_state=42} – ziarno losowości, zapewniające odtwarzalność podziałów danych w walidacji krzyżowej i podziale treningowym/testowym.
	\item \textbf{threshold=0.5} – próg prawdopodobieństwa, powyżej którego predykcje są klasyfikowane jako klasa pozytywna (1), stosowany w konwersji prawdopodobieństw na etykiety binarne.
	\item \textbf{zero\_division=0} – wartość domyślna dla obsługi dzielenia przez zero w metrykach precyzji i recall, zapobiegająca błędom przy braku predykcji jednej z klas.
	\item \textbf{flatten=True} – parametr określający, czy predykcje z modeli Keras powinny być spłaszczone do jednowymiarowej tablicy, ułatwiając dalsze przetwarzanie.
\end{itemize}

Drugim krokiem jest analiza i wizualizacja wyników, realizowana przez skrypt generujący wykresy porównawcze. Skrypt ten przetwarza wyniki z folderów z wynikami testów, obejmujących uruchomienia metod od 1 do 20 oraz dodatkowe pliki porównawcze. Wyniki ze wszystkich plików są łączone w jedną tabelę danych z dodaną kolumną \texttt{Metoda} wskazującą numer metody, których posiadamy 20. Dla metryki \texttt{Execution Time (s)} usuwane są wartości odstające. Następnie generowane są wykresy słupkowe dla wszystkich wspominanych metryk dla danych metod autorskich oznaczonych na czerwono i naukowych na niebiesko. Dodatkowo, dla wszystkich uruchomień generowany jest wykres krzywych Precision-Recall na podstawie danych z plików \texttt{runX\_pr\_curve.csv}.